\section{Metodología}

\subsection{Diseño centrado en el usuario}
El desarrollo del Portal Agro Comercial del Huila se basa en un enfoque centrado en el usuario, priorizando la accesibilidad y usabilidad para agricultores sin experiencia previa en tecnología. Se realizó un análisis de necesidades mediante entrevistas con productores rurales y observaciones en veredas del Huila, identificando barreras tecnológicas y culturales.

Hablemos claro: el dinero sigue siendo el mayor obstáculo para que la tecnología llegue a las fincas colombianas. La realidad es que muchísimos productores trabajan con márgenes de ganancia muy ajustados. Esto significa que casi no les sobra nada después de cubrir los gastos. En estas condiciones, ¿en qué se concentra el campesino? En lo vital, lo que garantiza la cosecha y la comida: Semillas y fertilizantes, Transporte de la cosecha, Y, por supuesto, el sostenimiento de la familia. Ante estas prioridades, ¿quién puede pensar en comprar un computador, pagar un buen plan de internet, o invertir en maquinaria avanzada? La tecnología, por necesaria que sea, queda relegada a un segundo plano, simplemente porque no hay presupuesto para ello. Por esta razón, cualquier plataforma digital que busque ser implementada debe ser económicamente accesible, con bajo costo operativo y alto retorno en valor comercial. Una política tecnológica para el agro requiere modelos de financiamiento, subsidios y alianzas entre Estado, academia y empresa para reducir el gasto del productor y ofrecer opciones escalables.

Aplicación al proyecto: El portal incluirá un plan totalmente gratuito de uso básico y suscripciones opcionales mejoradas para grandes productores. También se gestionarán convenios con alcaldías y asociaciones campesinas para brindar acceso a internet y equipos compartidos.

\subsection{Estrategias de adopción tecnológica}
Se implementó un diseño amigable con interfaces intuitivas, botones grandes, colores accesibles y lenguaje natural, evitando términos técnicos. Se incorporaron mecanismos de ayuda integrada, incluyendo tutoriales en video, chat de asistencia remota y capacitaciones presenciales en las comunidades.

La realidad de internet en el campo del Huila es como un mosaico: no es igual para todos. Hay veredas donde la señal 4G vuela, y otras donde apenas hay un poquito de cobertura, o solo funciona con un operador específico, ¡y a veces se cae! Esta "brecha digital" nos frena en muchos aspectos: Limita la comunicación básica, Hace casi imposible el comercio virtual, Dificulta acceder a cursos y educación técnica, Impide llevar un buen registro de los datos de la cosecha. Para que nuestros campesinos puedan entrar de lleno en la economía digital, el acceso a internet debe dejar de ser un lujo y pasar a ser un derecho básico, igual que el agua potable o la energía eléctrica.

Nuestra Solución para el Proyecto: Somos conscientes de que la señal falla, ¡y eso no puede detener el trabajo! Modo "Sin Señal": Nuestro portal tendrá una función que le permitirá al campesino registrar sus productos incluso sin internet. En cuanto la señal regrese, los datos se sincronizarán automáticamente. ¡El trabajo nunca se pierde! Impulso Comunitario: Propondremos a la Gobernación y a las Alcaldías un plan para instalar puntos de Wi-Fi comunitario en las veredas, especialmente para los campesinos que usen nuestro sistema. Así, aseguramos que la operación sea continua y reforzamos la inclusión digital en el campo.

\subsection{Funcionalidades offline y sincronización}
Para abordar la conectividad intermitente, se desarrolló un modo offline que permite registrar productos sin conexión a internet, con sincronización automática cuando se restablece la señal. Esta funcionalidad asegura la continuidad operativa en áreas rurales con cobertura limitada.

\subsection{Validación y evaluación}
Se realizaron pruebas piloto en veredas seleccionadas, evaluando la adopción mediante métricas de uso, retroalimentación de usuarios y tasas de retención. Se documentaron lecciones aprendidas para refinar el portal y garantizar su efectividad en la eliminación de intermediarios y promoción del comercio directo.